\lysh{16141_SOLUTION}{1}
ΛΥΣΗ

% TODO: TikZ σχήμα (μην το κατασκευάσεις)

$\alpha)$
\begin{enumerate}
    \item $(\vec{AB}, \vec{A\Gamma}) = 60^\circ$.
    \item $(\vec{AM}, \vec{B\Gamma}) = 90^\circ$.
    \item $(\vec{AM}, \vec{\Gamma A}) = 180^\circ - \widehat{\Gamma AM} = 180^\circ - 30^\circ = 150^\circ$. Καθώς η διάμεσος προς τη βάση είναι και διχοτόμος της γωνίας της κορυφής.
    \item $(\vec{BM}, \vec{\Gamma M}) = 180^\circ$.
    \item $(\vec{\Gamma M}, \vec{\Gamma B}) = 0^\circ$.
\end{enumerate}

$\beta)$ Τα μέτρα των διανυσμάτων $|\vec{AB}|, |\vec{A\Gamma}|$ και $|\vec{B\Gamma}|$ είναι $10$ αφού το τρίγωνο είναι ισόπλευρο. Για το μέτρο του διανύσματος $|\vec{AM}|$ έχουμε από το πυθαγόρειο θεώρημα στο τρίγωνο $AM\Gamma$:
\[
AM^2 = A\Gamma^2 - M\Gamma^2 \Rightarrow AM^2 = 10^2 - 5^2 = 75 \Rightarrow AM = \sqrt{75} = 5\sqrt{3}.
\]
\begin{enumerate}
    \item[i.] $\vec{AM} \cdot \vec{B\Gamma} = |\vec{AM}| \cdot |\vec{B\Gamma}| \sigma\upsilon\nu (\vec{AM}, \vec{B\Gamma}) = 5\sqrt{3} \cdot 10 \cdot \sigma\upsilon\nu 90^\circ = 0$.
    \item[ii.] $\vec{AM} \cdot \vec{\Gamma A} = |\vec{AM}| \cdot |\vec{\Gamma A}| \sigma\upsilon\nu (\vec{AM}, \vec{\Gamma A}) = 5\sqrt{3} \cdot 10 \cdot \sigma\upsilon\nu 150^\circ = 5\sqrt{3} \cdot 10 \cdot \left(-\frac{\sqrt{3}}{2}\right) = -75$.
    \item[iii.] $\vec{\Gamma M} \cdot \vec{\Gamma B} = |\vec{\Gamma M}| \cdot |\vec{\Gamma B}| \sigma\upsilon\nu (\vec{\Gamma M}, \vec{\Gamma B}) = 5 \cdot 10 \cdot \sigma\upsilon\nu 0^\circ = 5 \cdot 10 \cdot 1 = 50$.
\end{enumerate}
\end{lysh}